\section{\new{Extended Results}}

\subsection{\new{Statistical confidence and stability under small evaluation sets}}

\new{All primary comparisons carry 95\% percentile bootstrap confidence intervals ($B=2000$ document-level resamples), and pairwise CLINES-versus-baseline differences were tested with a paired bootstrap and Benjamini--Hochberg false-discovery-rate correction (Supplementary Table~\ref{supp:detailed-metrics}). CLINES (o3-mini-medium) exceeded both backbone alternatives on code, assertion, value, and unit extraction, with absolute F1 gains spanning $+0.01$ to $+0.21$ versus GPT-4o and Llama-3.1-405B (smallest gap on 4CE Unit, largest on CORAL--Breast Code). The Code-field comparisons (largest effect sizes) and most assertion/value gains were significant at $p_{\mathrm{FDR}}<0.05$; the smallest per-field gaps (e.g.\ 4CE Unit at $+0.01$) fall within bootstrap noise and were not significant. On the date columns, differences among the strongest backbones were not statistically resolvable at the available sample sizes. A document-level resampling stability analysis quantified the uncertainty arising from the small annotated sets: on CORAL--Breast ($n=13$) the bootstrap standard deviation of F1 reaches $\approx 0.04$--0.05, so F1 differences below $\approx 0.10$ are not reliably distinguishable. We therefore rest the principal claims on the comparisons whose effect sizes exceed this threshold, and Figure~3E annotates per-bin note counts to discourage over-interpretation of sparsely populated bins.}

\subsection{\new{Comparison with single-prompt and deep-learning baselines}}

\new{To isolate the contribution of the multi-step architecture from that of the backbone model and the prompt style, we compared CLINES against two single-call baselines and two full-size clinical deep-learning encoders (Figure~4d). On code-level extraction, o3-mini single-prompt reached F1 0.50 / 0.43 / 0.36 (4CE / CORAL--Pancreas / CORAL--Breast) and GPT-4o chain-of-thought 0.40 / 0.34 / 0.27, versus CLINES 0.87 / 0.85 / 0.81---a gain of $+0.37$ to $+0.54$ F1 attributable to the architecture rather than to the backbone model or to chain-of-thought prompting. The deep-learning encoders (BERT-base clinical NER, 110M; GatorTron-base, 345M, derived from the UF/NVIDIA GatorTron encoder) were competitive on \emph{mention detection} (F1 0.875 / 0.880 / 0.888 and 0.840 / 0.847 / 0.876, respectively, versus CLINES 0.906 / 0.881 / 0.913) but emit entity spans only: they perform no UMLS normalization, assertion, value+unit, or date extraction, so their code-level and attribute F1 are undefined. Classical rule- and lexicon-based systems (cTAKES, MetaMap) trailed substantially on entities; among them only cTAKES produced assertion labels, and neither produced value+unit or date outputs.}

\subsection{\new{Cost, latency, and computational efficiency}}

\new{Per-note cost varied by an order of magnitude across deployment options (Figure~4b). The reference CLINES (o3-mini-medium) configuration cost approximately \$0.21 per note in hosted-API charges, versus \$0.97 per note for the o3-mini single-prompt reasoning configuration, whose long reasoning traces inflate completion-token counts even at the same backbone model. The open-weight Llama-3.1-405B (FP8, 8$\times$H100) required $\approx 0.4$ GPU-hours per note. Because each note triggers several LLM invocations per chunk, the per-note call count exceeds the four architectural stages; the full per-step invocation, token, time, and cost breakdown is provided in Figure~4b.}

\subsection{\new{Component ablation}}

\new{Disabling each pipeline component in turn (Figure~4c) showed that SapBERT-based UMLS normalization is the load-bearing component: removing it reduced code-level F1 by 0.447 (GPT-4o) and 0.414 (GPT-4o-mini) on average. Cross-chunk reconciliation contributed $-0.078$ (GPT-4o) and $-0.221$ (GPT-4o-mini), semantic chunking $-0.021$ / $-0.031$, and the date module $-0.003$ / $-0.008$. Notably, reconciliation mattered more for the cheaper backbone, indicating that the deterministic scaffolding partly compensates for a weaker model---an argument in favor of the architecture precisely in cost-constrained settings. The ablation used two representative notes per dataset and reports point estimates.}

\subsection{\new{Performance on MIMIC-III}}

\new{On the MIMIC-III subset (24 expert-annotated paragraphs), CLINES reached entity F1 0.69, assertion F1 0.93, and value+unit F1 0.90. The lower entity F1 relative to 4CE and CORAL reflects (i)~the single-annotator-per-paragraph protocol used for MIMIC-III, which raises gold-standard variability and which we accordingly interpret as agreement-with-annotator rather than against a consensus reference, and (ii)~the heterogeneous, abbreviation-dense nature of critical-care notes, whose ICU-specific terminology is under-represented in the few-shot examples. Because MIMIC-III uses a single-annotator-per-paragraph protocol, we report it descriptively and confine the bootstrap and paired-significance analyses to the three consensus-annotated corpora.}

\subsection{\new{Positioning versus contemporary production and research systems}}

\new{Contemporary production and research systems address overlapping goals: production clinical-NLP platforms such as John Snow Labs Healthcare NLP / Spark NLP~\cite{JohnSnowLabs_HealthcareNLP} provide broad clinical NER, assertion, multi-terminology resolution, and OMOP output within hybrid LLM pipelines; CLEAR~\cite{Lopez_2025_CLEAR} is a closely related retrieval-augmented LLM pipeline for clinical extraction; and a recent multi-model evaluation~\cite{Ntinopoulos_2025} benchmarked output consistency across many LLMs. CLINES is positioned not as a competitor to mature production suites but as a transparent, fully-specified, model-agnostic open reference implementation that prioritizes reproducibility of its prompt suite and an auditable intermediate representation. A fair matched comparison (identical preprocessing, gold standard, and span-matching protocol) is a substantial undertaking left to future work.}
